\section{Oriented radial recovery and symplectic transport}
\subsection{Recovery of Planck's action with orientation preserved}

With the labelled semiaxes of the ellipse, define
\(R_\star=\sqrt{b_S/a_S}\). Then
\(q=R_\star^4=e^{-2\eta_{\rm el}}=(1-\xi)/(1+\xi)\),
where \(\xi=C^*/A\) and \(\eta_{\rm el}=\operatorname{artanh}\xi\).
This radial ratio retains the markings of both axes.

\begin{theorem}[Radial recovery map for the return section]
\label{rec:thm:accion}
In the positive chart, if \(q=R_\star^4\), then
\begin{align}
 \eta_{\rm ret}&=\alpha\frac{499-501q}{1+q},\label{rec:eq:eta-radio}\\
 \hbar_{\rm ret}&=\alpha\frac{499-501q}{1+q}\,10^{-34}\mathcal U_S,
 \qquad\hbar_{\rm pre}=H_5\,10^{-34}\mathcal U_S,\label{rec:eq:planck-radio}\\
 R_{\rm act}&=\frac{\hbar_{\rm pre}}{\hbar_{\rm ret}}
 =\frac{H_5(1+q)}{\alpha(499-501q)}.\label{rec:eq:ratio-radio}
\end{align}
The branch \(\eta_{\rm ret}>0\) is equivalent to \(q<499/501\), within \(q>0\).
\end{theorem}
\begin{proof}
The angular identities give \(\eta_{\rm ret}=\alpha(500\xi-1)\). Substitute \(\xi=(1-q)/(1+q)\); the numerator is \(499-501q\). Both the denominator and \(\alpha\) are positive. The dimensional realization and the ratio are those of the action sections already constructed. The decade \(-34\) and the unit \(\mathcal U_S\) retain their preceding determinantal derivation.
\end{proof}

This formula makes the Planck section within the counter-angle explicit. In terms of the action quantities, the angular relation is
\[
 C^*=2\left(\frac{\hbar_{\rm ret}}{10^{-34}\mathcal U_S}+\alpha\right).
\]
Division by the dimensional basis extracts the dimensionless coefficient; it does not add an action quantity to a dimensionless constant.

For the angular orientation \(\varepsilon\in\{-1,+1\}\), define
\[
 C^*_{\varepsilon}=2\varepsilon(\eta_{\rm ret}+\alpha),\quad
 \xi_\varepsilon=C^*_{\varepsilon}/A,\quad
 q_\varepsilon=\frac{1-\xi_\varepsilon}{1+\xi_\varepsilon}>0.
\]
The recovery map on both charts is
\begin{equation}
 \eta_{\rm ret}=\alpha\left(500\varepsilon
        \frac{1-q_\varepsilon}{1+q_\varepsilon}-1\right).
 \label{rec:eq:eta-orientada}
\end{equation}

\begin{corollary}[Invariance of the conjugate recovery map]
The transformation \((\varepsilon,q)\mapsto(-\varepsilon,q^{-1})\) preserves \(\eta_{\rm ret}\), \(\hbar_{\rm ret}\), and \(R_{\rm act}\).
\end{corollary}
\begin{proof}
The ratio \((1-q)/(1+q)\) changes sign when \(q\) is inverted. Its product with \(\varepsilon\) remains invariant. Apply \eqref{rec:eq:eta-orientada}, followed by the action representations.
\end{proof}

If orientation is discarded and \(\varepsilon=+1\) is incorrectly held fixed, the same expression gives \(\eta_{\rm ret}(q^{-1})=-\eta_{\rm ret}(q)-2\alpha\). The register therefore distinguishes two algebraically different operations. Nor is the angular orientation \(\varepsilon\) automatically identified with the symplectic sign \(\sigma\): exchanging the axes is antisymplectic, whereas the rotation that exchanges them preserves the area form.

\subsection{Intertwining phase and anisotropy}

Let
\[
 S_\eta=\operatorname{diag}(e^{\eta/2},e^{-\eta/2}),\quad
 J_0=\begin{pmatrix}0&-1\\1&0\end{pmatrix},\quad
 J_\eta=S_\eta J_0S_\eta^{-1}
       =\begin{pmatrix}0&-e^\eta\\e^{-\eta}&0\end{pmatrix}.
\]
Then \(\det S_\eta=1\), \(J_\eta^2=-I\), and
\(C_\eta(\theta)=e^{\theta J_\eta}=S_\eta e^{\theta J_0}S_\eta^{-1}\).
The law \(S_{\eta+\delta}=S_\delta S_\eta\) proves the exact intertwining identity
\begin{equation}
 C_{\eta+\delta}(\theta)S_\delta=S_\delta C_\eta(\theta).
 \label{rec:eq:entrelazamiento}
\end{equation}
The transformation preserves \(dQ\wedge dP\) and carries the phase evolution to the new ellipse. With the normalization of \eqref{eq:ii-elipse-area-fase}, the swept action is \(\hbar\theta\), or \(\sigma\hbar\theta\) on the oriented sheet. The phase \(\theta\), the anisotropy \(\eta_{\rm el}\), and the parameter \(t\) of the flow of \(K\) thus retain distinct functions, related by explicit operators.
